% Native publication source. Preserve the accompanying semantic table and caption.
% Requires \usepackage{tikz}. Add tagged figure metadata through the final template.
\begin{figure}[htbp]
\centering
\begin{tikzpicture}[alt={Seven-segment positions labelled a through g and dp, corresponding to bits zero through seven.},x=1cm,y=1cm,line width=4pt,line cap=round]
\draw (0.2,4)--(1.8,4); \draw (2,3.8)--(2,2.2);
\draw (2,1.8)--(2,0.2); \draw (0.2,0)--(1.8,0);
\draw (0,0.2)--(0,1.8); \draw (0,2.2)--(0,3.8);
\draw (0.2,2)--(1.8,2); \fill (2.5,0) circle (0.09);
\node[above] at (1,4.1) {a: bit 0};
\node[right] at (2.1,3) {b: bit 1};
\node[right] at (2.1,1) {c: bit 2};
\node[below] at (1,-0.1) {d: bit 3};
\node[left] at (-0.1,1) {e: bit 4};
\node[left] at (-0.1,3) {f: bit 5};
\node[above] at (1,2.1) {g: bit 6};
\node[below right] at (2.5,-0.1) {dp: bit 7};
\end{tikzpicture}
\caption{Segment positions and their mask bits. A set bit lights the named segment; bit 7 controls the decimal point.}
\label{fig:counter-segment-map}
\end{figure}
